\anglais

\chapter*{General Introduction}
\label{chap:general_introduction}

Modern society depends on technologies that are exposed, directly or indirectly, to the activity of the Sun. Satellites, radio communication, navigation systems, aviation, and electrical power networks can all be affected by severe space-weather events. Among these events, solar flares are sudden releases of magnetic energy in the solar atmosphere. They emit intense radiation across a broad range of wavelengths and typically occur in active regions where the solar magnetic field is complex and rapidly evolving. Forecasting these events is therefore important not only as a scientific problem, but also as a practical step toward reducing the impact of space weather on technological infrastructure.

Solar flares are classified according to their peak X-ray flux into different intensity classes. Moderate M-class flares can already affect radio communication and cause minor satellite anomalies, while the most intense X-class flares can lead to severe disruptions, including satellite damage, navigation errors, and geomagnetically induced currents that may affect power grids on Earth.

These events are also increasingly rare with intensity: C-class flares are relatively common, M-class flares occur much less frequently, and X-class flares are rare extreme events. This combination of high impact and low frequency makes their prediction particularly challenging and operationally important.

Solar flare prediction is difficult because the relevant information is both spatial and temporal. Active regions can be observed through magnetograms, images that describe the magnetic field at the solar surface. These images contain polarity patterns, gradients, and large-scale magnetic organization that may be related to future eruptions. However, a flare is not simply an object visible in an image. It is a future event, and two active regions that look similar at one instant may evolve differently over the following hours. This makes the problem different from standard image classification: the model must learn visual precursors of future solar activity rather than only recognize present structures.

The task is also strongly imbalanced. Moderate flares are more common than intense ones, while stronger M- and X-class flares (see Chapter~\ref{chap:introduction} for a detailed definition of flare classes) are rare. This imbalance affects both training and evaluation. A forecasting model must detect meaningful positive cases without being rewarded for simply predicting the majority negative class. For this reason, solar flare prediction studies often rely on skill scores and ranking metrics, such as the True Skill Statistic, ROC-AUC, and precision-recall AUC, rather than accuracy alone.

Deep learning provides a promising framework for this problem because it can learn representations directly from magnetogram images and sequences. Convolutional models, vision transformers, and spatio-temporal architectures can capture different aspects of active-region structure. In particular, temporal models can use several observations of the same active region to represent how magnetic patterns change before a flare. At the same time, deep learning pipelines require careful design choices: the input resolution, the position of the active region on the solar disk, the cadence between frames, and the number of frames can all influence performance, computation, and the amount of usable data.

This thesis investigates input-design choices for solar flare prediction from magnetogram data, using strong pretrained deep-learning models within a controlled fine-tuning framework to isolate the effect of input representation on performance. The objective is to determine which design decisions preserve predictive skill and to provide practical guidance for building efficient computer-vision pipelines for space-weather forecasting.

The thesis is organized as follows. Chapter~\ref{chap:introduction} introduces the scientific and computational background of solar flare prediction, including magnetogram data, flare labels, evaluation metrics, dataset properties, and the deep-learning architectures used in this work. Chapter~\ref{chap:article1} presents the main article-style study, which evaluates the effect of spatial resolution, solar-disk position, temporal cadence, and number of input frames on flare prediction performance. The thesis concludes by summarizing the main findings and outlining directions for future work. The appendices, beginning with Appendix~\ref{app:supplementary_empirical_analysis}, provide additional empirical material supporting the analyses presented in Chapter~\ref{chap:article1}.
